\title{xLSTM: Extended Long Short-Term Memory}

\begin{abstract}
The introduction of the constant error carousel and gating mechanisms in the 1990s established Long Short-Term Memory (LSTM) as foundational components for handling long-term dependencies in sequential data. Over subsequent decades, LSTMs played a pivotal role in advancing deep learning, most notably serving as the core architecture behind the earliest Large Language Models (LLMs). The rise of the Transformer framework, built upon highly parallelizable self-attention mechanisms, marked a paradigm shift, eventually eclipsing LSTMs at scale. Motivated by these developments, we ask: What are the limits of language modeling performance when we scale LSTMs to billions of parameters, draw upon cutting-edge LLM advancements, and directly address the classic challenges faced by LSTM architectures? To this end, we propose exponential gating enhanced by normalization and stabilization strategies. In addition, we re-engineer the LSTM’s memory structure, resulting in: (i) sLSTM, which features scalar memory, scalar update steps, and a novel mixing mechanism, and (ii) mLSTM, enabling full parallelization by leveraging a matrix-based memory with a covariance-based update rule. By embedding these enhanced LSTM variants within residual block frameworks, we introduce xLSTM blocks, which are subsequently assembled in a residual fashion to form the xLSTM model family. Our approach—encompassing exponential gating and restructured memory—substantially improves the scalability and competitive performance of xLSTM models, demonstrating advantages over contemporary Transformers and State Space Models in both efficiency and effectiveness.\\
Code available at: \url{https://github.com/NX-AI/xlstm}
\end{abstract}

\section{Introduction}
\label{sec:intro}

Long Short-Term Memory (LSTM) concepts~\citep{Hochreiter:91,Hochreiter:97nips,Hochreiter:97}, specifically the mechanisms of constant error carousel and gating, were developed to address the vanishing gradient issue present in recurrent neural networks~\citep{Hochreiter:91,Hochreiter:00book}:

\begin{align}
\label{eq:lstm_idea}
  c_t \ = \ f_t \ c_{t-1} \ + \ 
          i_t \  z_t \ , \quad  
          h_t \ = \ o_t \ \psi({c_t}) \ . 
\end{align}

The constant error carousel operates by adding cell inputs $z_t$ to the previous cell state $c_{t-1}$ (indicated in green), with this combination regulated through sigmoid gates (shown in blue). The input gate $i_t$ and forget gate $f_t$ govern how the cell state is updated, whereas the output gate $o_t$ determines the memory cell’s output, specifically the hidden state $h_t$. Prior to being gated for output, the cell state undergoes normalization or squashing via $\psi$, after which the output gating generates the hidden state.

LSTMs have achieved considerable success across multiple fields \citep{Hochreiter:01,Hochreiter:07,Schmidhuber:15}, dominating text generation tasks prior to the emergence of Transformers in 2017~\citep{Vaswani:17}. Their capabilities have been validated through a wide range of sequence-processing challenges, such as text generation~\citep{Graves:13,Karpathy:15blog}, handwriting synthesis~\citep{Graves:13}, sequence-to-sequence translation~\citep{Sutskever:14nips}, program evaluation~\citep{Zaremba:14arxiv}, captioning images~\citep{Karpathy:15,Hossain:19}, producing source code~\citep{Karpathy:15blog}, rainfall-runoff prediction~\citep{Kratzert:18,Kratzert:19}, and constructing hydrological models for flood warnings~\citep{Nearing:24}. Within reinforcement learning, LSTMs remain leading sequence models, exemplified by their use in systems such as AlphaStar for StarCraft II~\citep{Vinyals:17short}, OpenAI Five for Dota 2~\citep{OpenAI:19}, and the magnetic controller models for nuclear fusion~\citep{Degrave:22short}. One of the key strengths of LSTMs is their ability to form high-level abstractions—efficiently extracting and preserving semantic content in their cell states~\citep{Karpathy:15blog}. This skill has been observed through the identification of number and syntax neurons~\citep{Lakretz:19}, linguistic-specific neurons~\citep{Bau:19}, and sentiment neurons~\citep{Radford:17arxiv}. LSTMs remain prominent in contemporary, mission-critical applications~\citep{Degrave:22short,Nearing:24}, further attesting to their enduring relevance and robustness.

Although LSTMs have achieved remarkable progress, they exhibit three primary drawbacks: (i) Inability to revise storage decisions. To illustrate, consider the {\it Nearest Neighbor Search} task (refer also to Appendix~\ref{sec:appNearestNeighborSearch}): Given a reference vector, one must examine a sequence in order to locate the most similar vector and then return the value associated with it at the end of the sequence. As depicted in the left panel of Figure~\ref{fig:lstmProblems}, this results in mean squared error that reveals a significant weakness—LSTM fails to update a previously stored value if a closer vector appears later in the sequence. By contrast, our xLSTM overcomes this constraint through the application of exponential gating. (ii) Restricted storage capacity—information is necessarily compacted into scalar cell states. This drawback is demonstrated through {\it Rare Token Prediction}. In the right panel of Figure~\ref{fig:lstmProblems}, we report the token prediction perplexity on the Wikitext-103 dataset~\citep{Merity:17}, separating results by token frequency. The LSTM exhibits degraded performance for infrequent tokens due to these limited storage abilities, whereas the matrix-based memory of xLSTM effectively addresses this issue. (iii) Inherent challenges with parallelization caused by memory mixing, specifically the recurrent connections that propagate hidden state information across time steps and thus require inherently sequential processing.

The shortcomings inherent to LSTM architectures have contributed to the development and widespread adoption of Transformers~\citep{Vaswani:17} in the field of language modeling. An open question remains: what level of performance could be attained in language modeling tasks if one addresses these constraints and expands LSTMs to match the scale of contemporary Large Language Models?

\section{Extended Long Short-Term Memory}
\label{sec:xLSTM}

To address the constraints of standard LSTMs, the Extended Long Short-Term Memory (xLSTM) framework proposes two significant changes to the conceptual formulation of LSTM in Equation~\eqref{eq:lstm_idea}. These changes—namely, the adoption of exponential gating mechanisms and the incorporation of new memory structures—result in two variants: (i) sLSTM (see Section~\ref{sec:sLSTM}), which utilizes scalar memory, scalar updates, and a memory mixing procedure; and (ii) mLSTM (see Section~\ref{sec:mLSTM}), which is characterized by matrix-valued memory and covariance-based (outer product) updates that are fully parallelizable. Both the sLSTM and mLSTM models extend the classical LSTM by leveraging exponential gating. Notably, to achieve parallel computation, the mLSTM omits the memory mixing step—specifically, the hidden-to-hidden recurrent connections. Both architectures, mLSTM and sLSTM, permit the use of multiple memory cells, where sLSTM specifically implements mixing across these cells. Furthermore, the sLSTM architecture can support multiple heads; in such cases, memory mixing is absent between different heads but remains present among cells within a head. The combination of multiple heads in sLSTM with exponential gating leads to a novel form of memory interaction. For mLSTM, however, there is no distinction between the introduction of multiple heads and that of multiple cells—both serve the same function.

The incorporation of these novel LSTM variants within residual block modules leads to the creation of xLSTM blocks (refer to Section~\ref{sec:xLSTMblock}). When these xLSTM blocks are stacked residually, they form xLSTM architectures, as detailed in Section~\ref{sec:xLSTMarchitecure}. Figure~\ref{fig:xlstm_sketch} illustrates the structure of the xLSTM architecture along with its constituent parts.

\subsection{Review of the Long Short-Term Memory}
\label{sec:orgLSTM}

The foundational concept of LSTM~\citep{Hochreiter:91,Hochreiter:97nips,Hochreiter:97} established the scalar memory cell as its core component for information processing and retention. This architecture effectively addresses the problem of vanishing gradients~\citep{Hochreiter:91,Hochreiter:00book} by implementing the constant error carousel mechanism, which is reflected in the update of the cell state. Within the memory cell, three types of gates regulate information flow: the input gate, the output gate, and the forget gate, the latter of which was later incorporated by \citet{Gers:00}. The dynamics governing the update of the LSTM memory cell at a particular time step $t$ are described as follows:

\begin{align}
c_t \ &= \  f_t \ c_{t-1} \ + \ i_t \ z_t &  & &\text{cell state} \\
h_t  \ &= \ o_t \ \tilde{h}_t \ , & \tilde{h}_t \ &= \ \psi \left( c_t\right)
  &\text{hidden state} \\
z_t \ &= \ \varphi \left( \tilde{z}_t \right) \ , 
  &\tilde{z}_t \ &=  \ \mathbf{w}^\top_{z} \ \mathbf{x}_t \ + \
  r_{z}  h_{t-1} \ + \  b_{z} \ \
  &\text{cell input} \\
i_t \ &= \ \sigma \left( \tilde{i}_t  \right) \ , 
  &\tilde{i}_t \ &= \ \mathbf{w}^\top_{i} \ \mathbf{x}_t \ + \
  r_{i}  \ h_{t-1} \ + \  b_{i} \ \
  &\text{input gate} \\
f_t \ &= \ \sigma \left(  \tilde{f}_t \right) \ , 
  &\tilde{f}_t \ &= \ \mathbf{w}^\top_{f} \ \mathbf{x}_t  \ + \
  r_{f}  \ h_{t-1} \ + \  b_{f} \ \
  &\text{forget gate} \\
o_t \ &= \ \sigma \left( \tilde{o}_t \right) \ , 
  &\tilde{o}_t  \ &= \ \mathbf{w}^\top_{o} \ \mathbf{x}_t \ + \
  r_{o}  \ h_{t-1} \ + \  b_{o} \ \
  &\text{output gate} 
\end{align}

The vectors $\mathbf{w}_{z}$, $\mathbf{w}_{i}$, $\mathbf{w}_{f}$, and $\mathbf{w}_{o}$ represent the input weights linking the input $\mathbf{x}_t$ to the cell input, the input gate, the forget gate, and the output gate, respectively. Similarly, the weights $r_{z}$, $r_{i}$, $r_{f}$, and $r_{o}$ serve as recurrent connections from the previous hidden state $h_{t-1}$ to each of these components. The associated biases are denoted as $b_{z}$, $b_{i}$, $b_{f}$, and $b_{o}$. The activation functions for the cell input and hidden state are $\varphi$ and $\psi$, which are often chosen to be $\tanh$. The activation function $\psi$ is applied to the cell state to constrain its values, preventing it from becoming unbounded. The gating mechanisms employ the sigmoid activation, $\sigma \left( x \right)= 1/(1+\exp(-x))$. In subsequent designs, collections of individual memory cells, $\mathbf{c}_t \in \mathbb{R}^{d}$, were combined into a vector, making it possible to utilize recurrent weight matrices $\mathbf{R} \in \mathbb{R}^{d \times d}$ that integrate information across memory cells~\citep{Greff:15}. Additional technical discussion can be found in Appendix~\ref{sec:apporgLSTM}. Empirical ablation analyses have indicated that every part of the memory cell is essential for effective functioning~\citep{Greff:15}.

\subsection{sLSTM}
\label{sec:sLSTM}

To enable LSTMs to reconsider their memory storage choices, we propose the addition of exponential gates (depicted in red) alongside mechanisms for normalization and stabilization. Specifically, the input and forget gates are allowed to employ exponential activation functions. For the purpose of normalization, we define a normalizer state that accumulates the sum of the input gate multiplied by the subsequent forget gates over time.

The scalar sLSTM forward pass is formulated as:

\begin{align}
\label{eq:slstmforward}
c_t \ &= \  f_t \ c_{t-1} \ + \ i_t \ z_t & & 
 &\text{cell state} \\
n_t \ &= \  f_t \ n_{t-1} \ + \ i_t  & & 
 &\text{normalizer state} \\
h_t \ &= \ o_t \ \tilde{h}_t \ ,
  &\tilde{h}_t \ &= \ c_t / n_t
&\text{hidden state} \\
z_t \ &= \ \varphi \left( \tilde{z}_t \right) \ , 
  &\tilde{z}_t \ &=  \ \mathbf{w}^\top_{z} \ \mathbf{x}_t \ + \
  r_{z}  h_{t-1} \ + \  b_{z}
  &\text{cell input} \\
i_t \ &= \ \!\exp \left( \tilde{i}_t  \right) \ , 
  &\tilde{i}_t \ &= \ \mathbf{w}^\top_{i} \ \mathbf{x}_t \ + \
  r_{i}  \ h_{t-1} \ + \  b_{i}
  &\text{input gate} \\
f_t \ &= \ \sigma \left(  \tilde{f}_t \right) \ \text{OR} \
\!\exp \left(  \tilde{f}_t \right) \ ,
  &\tilde{f}_t \ &= \ \mathbf{w}^\top_{f} \ \mathbf{x}_t  \ + \
  r_{f}  \ h_{t-1} \ + \  b_{f}
  &\text{forget gate} \\
o_t \ &= \ \sigma \left( \tilde{o}_t \right) \ , 
  &\tilde{o}_t  \ &= \ \mathbf{w}^\top_{o} \ \mathbf{x}_t \ + \
  r_{o}  \ h_{t-1} \ + \  b_{o}
  &\text{output gate} 
\end{align}

The gating mechanisms from standard LSTM architectures—namely those dependent on inputs and/or hidden states, together with the associated bias term—are adapted for application in these new models. Given that exponential activation functions may result in excessively large outputs, leading to potential overflow issues, we incorporate an extra state, denoted as \mbox{$m_t$~\citep{Milakov:18arxiv}}, to provide additional stability for the gates:

\begin{align}
\label{eq:slstmstabil}
m_t \ &= \ \max \left( \log ( f_t ) + m_{t-1} , \log ( i_t ) \right) &\text{stabilizer state} \\
i'_t \ &= \ \!\exp \left( \log \left ( i_t \right) - m_t \right) \ = \ \!\exp \left( \tilde{i}_t - m_t \right) \ \, 
  &\text{stabil. input gate} \\
f'_t \ &= \ \!\exp \left( \log \left( f_t \right) + m_{t-1} - m_t \right) \ \, 
  &\text{stabil. forget gate}
\end{align}

As demonstrated in Appendix~\ref{sec:appsLSTM}, substituting $f_t$ with $f'_t$ and $i_t$ with $i'_t$ during the forward computation leaves both the overall network output and the gradients of the loss function with respect to the parameters unaffected.

\paragraph{New Memory Mixing.} Similar to standard LSTMs (refer to Appendix~\ref{sec:appsLSTM}), sLSTM accommodates several memory cells. The use of multiple memory cells allows for memory mixing, facilitated by recurrent connections $\mathbf{R}_{\mathbf{z}}$, $\mathbf{R}_{\mathbf{i}}$, $\mathbf{R}_{\mathbf{f}}$, and $\mathbf{R}_{\mathbf{o}}$, which link the hidden state $\mathbf{h}$ with the memory cell input $\mathbf{z}$ and the corresponding gates $\mathbf{i}$, $\mathbf{f}$, and $\mathbf{o}$. A distinctive feature in this approach is the integration of exponential gating in the memory mixing process. The enhanced sLSTM supports multiple heads, where memory mixing is performed within a head but not between different heads. Incorporating heads in sLSTM, alongside the mechanism of exponential gating, thus introduces a novel framework for memory mixing.

\subsection{mLSTM}
\label{sec:mLSTM}

In order to improve the storage capability of LSTMs, we substitute the scalar memory cell $c \in \mathbb{R}$ with a matrix memory cell $\mathbf{C} \in \mathbb{R}^{d \times d}$. This modification enables information retrieval through matrix multiplication operations. Specifically, at a given time step $t$, our goal is to memorize a tuple consisting of a key vector $\mathbf{k}_t \in \mathbb{R}^d$ and an associated value vector $\mathbf{v}_t \in \mathbb{R}^d$, following the nomenclature of the Transformer architecture. Subsequently, at time $t+\tau$, we wish to recover the original value $\mathbf{v}_t$ using a query vector $\mathbf{q}_{t+\tau} \in \mathbb{R}^d$. This framework corresponds to the classical setup found in Bidirectional Associative Memories (BAMs) \citep{Kohonen:72,Anderson:72,Nakano:72,Anderson:77}.

The mechanism for updating the memory with a new key-value pair employs the covariance update rule~\citep{Sejnowski:77,Dayan:91}, which can be described as follows:

\begin{align}
\mathbf{C}_t \ &= \ \mathbf{C}_{t-1} \ + \ \mathbf{v}_{t} \ \mathbf{k}_t^\top \ .
\end{align}

We consider that inputs undergo layer normalization prior to their projection onto keys and values, ensuring these projections have zero mean. The covariance update mechanism has been shown to be optimal~\citep{Dayan:91} in terms of maximizing the separability of retrieved binary vectors, an outcome which coincides with achieving the highest possible signal-to-noise ratio. Enhancements in separability can be realized if retrieval is confined to pairwise interactions, albeit at the expense of quadratic computational complexity, similar to attention mechanisms~\citep{Krotov:16,Krotov:17,Ramsauer:21}. This covariance update process aligns with the Fast Weight Programmers paradigm~\citep{Schmidhuber:92ncfastweights,Schlag:21}, where subsequent developments have introduced a fixed decay rate for scaling $\mathbf{C}_{t-1}$ and a constant learning rate for the term 
$\mathbf{v}_{t} \mathbf{k}_t ^\top$~\citep{Ba:16}. Following this line of work, we incorporate the covariance update principle within the LSTM architecture: here, the forget gate modulates the decay rate, the input gate controls the learning rate, and the output gate adjusts the scale of the vector being retrieved.

Within this matrix memory framework, the normalizer state is computed as a weighted aggregation of key vectors, where each key's contribution is determined by the input gate and the sequence of subsequent forget gates. The normalizer state thus encapsulates the cumulative influence of the gating mechanisms. To address the possibility that the dot product between the query and the normalizer state approaches zero, we take the absolute value of this dot product and constrain it to be no less than a predefined threshold (commonly set to 1.0), following the methodology outlined in~\citep{Sun:23arxiv}. The mLSTM forward computation is:

\begin{align}
\label{eq:mlstm_recurrent_begin}
\mathbf{C}_t \ &= \  f_t \ \mathbf{C}_{t-1} \ + \ 
  i_t \ \mathbf{v}_t \ \mathbf{k}_t^\top & &  &\text{cell state} \\
\mathbf{n}_t \ &= \  f_t \ \mathbf{n}_{t-1} \ + \ 
  i_t \ \mathbf{k}_t & &  &\text{normalizer state} \\
\mathbf{h}_t  \ &= \ \mathbf{o}_t \ \odot \ \tilde{\mathbf{h}}_t
\ , \qquad \qquad 
\tilde{\mathbf{h}}_t \ = \   \mathbf{C}_t \mathbf{q}_t \ / \ 
  \max \left\{ |\mathbf{n}_t^\top \mathbf{q}_t|, 1 \right\} 
   &&&\text{hidden state} \\
\mathbf{q}_t \ &= \ \mathbf{W}_q \ \mathbf{x}_t \ + \ \mathbf{b}_q  & &  &\text{query input} \\
\mathbf{k}_t \ &= \ \frac{1}{\sqrt{d}} \mathbf{W}_k \ \mathbf{x}_t \ + \ \mathbf{b}_k  & &  &\text{key input} \\
\mathbf{v}_t \ &= \ \mathbf{W}_v \ \mathbf{x}_t \ + \ \mathbf{b}_v  & &  &\text{value input} \\
i_t \ &= \ \!\exp \!  \! \left( \tilde{i}_t  \right) \ , 
 \qquad \qquad \ \ \ \ \,
  \tilde{i}_t \ = \ \mathbf{w}^\top_{i} \ \mathbf{x}_t \ + \  b_{i} \ \ \ 
  &&&\text{input gate} \\
f_t \ &= \ \sigma \!  \left(  \tilde{f}_t \right) \ \text{OR} \ \!\exp \!  \! \left(  \tilde{f}_t \right) , \ \ \: \!
\tilde{f}_t \ = \ \mathbf{w}^\top_{f} \ \mathbf{x}_t  \ + \
  b_{f}
  &&& \text{forget gate} \\
\label{eq:mlstm_recurrent_end}
\mathbf{o}_t \ &= \ \sigma \left( \tilde{\mathbf{o}}_t \right) \ , \qquad \qquad \quad \ \ \ \,
  \tilde{\mathbf{o}}_t  \ = \ \mathbf{W}_{\mathbf{o}} \ \mathbf{x}_t \ + \
  \mathbf{b}_{\mathbf{o}}
  &&&\text{output gate} 
\end{align}

Similar to the original LSTM architecture, mLSTM supports several memory cells. In mLSTM, the use of multiple heads corresponds directly to multiple cells, due to the absence of any memory mixing. To ensure stability of the exponential gating mechanisms in mLSTM, we apply the same stabilization strategies utilized for sLSTM (refer to Equation~\ref{eq:slstmstabil}). Given that mLSTM does not incorporate memory mixing, its recurrence relations may be rewritten in a parallel form. Additional information can be found in Appendix~\ref{sec:appmLSTM}.

\subsection{xLSTM Architecture}
\label{sec:xLSTMarch}

\paragraph{xLSTM Blocks.}
\label{sec:xLSTMblock}
An xLSTM block is intended to create a non-linear, high-dimensional summary of preceding information, enhancing the distinction between different historical states or contexts. Achieving this separation is crucial for accurately forecasting the upcoming element in a sequence, such as the next token. Our design draws on Cover’s Theorem~\citep{Cover:65}, which indicates that patterns projected non-linearly into a higher-dimensional space become more susceptible to linear separation compared to their organization in the original space. We investigate two types of residual block structures: (i) The first uses a post up-projection, akin to the architecture of Transformers, where the non-linear processing of past inputs occurs in the base space, followed by a linear projection into a higher-dimensional space, application of a non-linear activation, and finally a linear projection back to the base space; this process is illustrated in the leftmost panel of Figure~\ref{fig:Backbones} as well as the third column of Figure~\ref{fig:xlstm_sketch}. For a comprehensive depiction, refer to Figure~\ref{fig:appsLSTM_detailed} in the appendix. (ii) The second employs a pre up-projection, as commonly found in State Space Models, which first linearly projects to a higher-dimensional representation,

The method captures historical information in a non-linear fashion within a high-dimensional space before applying a linear transformation to revert to the initial dimensionality. When an xLSTM block incorporates an sLSTM, we implement the up-projection block after this processing. Conversely, for xLSTM blocks that utilize an mLSTM, the up-projection block is placed beforehand, owing to the expanded memory in the higher-dimensional space. For additional clarification, consult the left panel of Figure~\ref{fig:Backbones}, the third column in Figure~\ref{fig:xlstm_sketch}, or see Figure~\ref{fig:appsLSTM_detailed} in the appendix.

\paragraph{xLSTM Architecture. }
\label{sec:xLSTMarchitecure}
The xLSTM model is designed by assembling its foundational components in a residual manner~\citep{Srivastava:15,He:16}. Our implementation follows the widely adopted pre-LayerNorm~\citep{Ba:16b} residual configuration, which is the standard approach in recent Large Language Models. Refer to the rightmost column in Figure~\ref{fig:xlstm_sketch}.

\subsection{Memory and Speed Considerations}

In contrast to Transformers, xLSTM networks feature linear computational requirements and maintain constant memory complexity regardless of sequence length. The compressive nature of xLSTM memory makes these networks particularly advantageous for industrial use cases and deployment on edge devices.

While the mLSTM's memory does not depend on any parameters, it is computationally intensive due to its reliance on a $d \times d$ matrix for memory storage and updates of the same dimensionality. This architecture sacrifices computational efficiency for increased memory capacity. However, since these calculations can be executed in parallel on GPUs, their impact on the actual elapsed time is minimal.

Although mLSTM can be executed in parallel similarly to FlashAttention~\citep{Dao:22,Dao:23} or GLA~\citep{Yang:23arxiv}, sLSTM lacks parallelizability because of its memory mixing involving hidden-to-hidden interactions. Nevertheless, we engineered an efficient CUDA version featuring optimizations down to the register layer in GPU memory, resulting in performance that is generally no more than twice as slow as mLSTM.

\section{Related Work}
\label{sec:related}

\paragraph{Linear Attention.} A number of approaches have been proposed to reduce the Transformer’s quadratic time complexity with respect to input sequence length, thereby achieving linear complexity. Synthesizer generates artificial attention weights independent of direct token-to-token computations~\citep{Tay:20arxiv}. Linformer implements self-attention through low-rank matrix projections, providing a linear approximation~\citep{Wang:20arxiv}. Linear Transformer recasts the attention process into a linear form~\citep{Katharopoulos:20}. Performer uses a method based on positive orthogonal random features to obtain a linear approximation of the softmax function in attention~\citep{Choromanski:21}. Additionally, alternative structures such as fast long convolutions have substituted attention in models like Structured Global Convolution (SGConv)~\citep{Li:22} and Hyena Hierarchy~\citep{Poli:23}.

\paragraph{State Space Models.} State Space Models (SSMs) have recently gained significant traction, primarily due to their linear scalability with respect to sequence length and their competitive results when compared to Transformers. Among the early contributions to this line of research was the Structured State Space sequence model (S4)~\citep{Gu:21}, which was subsequently followed by variants such as the Diagonal State Space (DSS) model~\citep{Gupta:22}, Gated State Space (GSS) models~\citep{Mehta:22}, the S5 model~\citep{Smith:22}, Bidirectional Gated SSM (BiGS)~\citep{Wang:22}, the H3 model~\citep{Fu:23}, and, most recently, Mamba~\citep{Gu:24arxiv}.

\paragraph{Recurrent Neural Networks.} Recurrent Neural Networks (RNNs) have been proposed as alternatives to Transformer architectures and attention mechanisms because of their ability to scale linearly with respect to context length. Notably, RNNs employing Deep Linear Recurrent Units (LRUs) have yielded strong performance in language modeling tasks~\citep{Orvieto:23, De:24arxiv}. Other architectures, such as the Hierarchically Gated Linear RNN (HGRN)~\citep{Qin:23} and its successor HGRN2~\citep{Qin:24arxiv}, have demonstrated similar promise. Among established RNN methodologies for large language modeling, the RWKV approach~\citep{Peng:23arxivshort,Peng:24arxivshort} is recognized for achieving results that are on par with those produced by Transformer-based models.

\paragraph{Gating.}
Gating, a fundamental mechanism originally introduced in LSTM architectures, has been revisited and reimagined across a range of recent methods. Variants of gating have been incorporated into HGRN~\citep{Qin:23}, HGRN2~\citep{Qin:24arxiv}, Gated Linear Attention (GLA)~\citep{Yang:23arxiv}, Gated State Space (GSS) models~\citep{Mehta:22}, Bidirectional Gated SSM (BiGS)~\citep{Wang:22}, Moving Average Equipped Gated Attention (MEGA)~\citep{Ma:22}, RWKV~\citep{Peng:23arxivshort}, and Mamba~\citep{Gu:24arxiv}.

\paragraph{Covariance Update Rule.} To improve the memory capacity, we augmented the mLSTM cell by incorporating a matrix memory governed by a covariance update rule. Comparable strategies that utilize a similar update framework can be found in Fast Weight Programmers \citep{Schmidhuber:92ncfastweights,Schlag:21}, RWKV-5 and RWKV-6~\citep{Peng:24arxivshort}, Retention~\citep{Sun:23arxiv}, Linear Transformer~\citep{Katharopoulos:20}, and HGRN2~\citep{Qin:24arxiv}.

\paragraph{Most Related.} The models bearing the strongest conceptual resemblance to xLSTM are Retention~\citep{Sun:23arxiv}, RWKV~\citep{Peng:23arxivshort,Peng:24arxivshort}, and HGRN2~\citep{Qin:24arxiv}. All of these architectures utilize mechanisms such as matrix memory and/or gating. Nonetheless, unlike the novel sLSTM, they lack support for memory mixing. The ability to mix memory is critical for addressing state tracking tasks, which renders LSTMs inherently more expressive than either State Space Models (SSMs) or Transformers~\citep{Merrill:24,Deletang:23}. Effective state tracking is vital for applications like program evaluation and following entities throughout extended narratives.

\paragraph{Residually Stacking Architectures.} In line with nearly all state-of-the-art deep learning frameworks, xLSTM architectures are assembled by stacking modular components with residual connections~\citep{Srivastava:15,He:16}.

Such an approach facilitated the successful development of deep convolutional neural networks~\citep{He:16} as well as the emergence of Transformer models~\citep{Vaswani:17}. The Transformer architecture, in particular, is foundational to Large Language Models (LLMs) such as GPT-3~\citep{Brown:20short}, ChatGPT~\citep{Schulman:22short}, GPT-4~\citep{Achiam:23gpt4short}, Megatron-LM~\citep{Shoeybi:19}, Gopher~\citep{Rae:21short}, ERNIE 3.0 Titan~\citep{Wang:21arxivshort}, GLaM~\citep{Du:21glamshort}, Chinese M6~\citep{Lin:21}, the multilingual AlexaTM 20B~\citep{Soltan:22}, OPT~\citep{Zhang:22opt}, Chinchilla~\citep{Hoffmann:22short}, BLOOM~\citep{Scao:22arxivshort}, GLM-130B~\citep{Zeng:22arxivshort}, LaMDA~\citep{Thoppilan:22short}, PaLM~\citep{Chowdhery:22short}, Llama~\citep{Touvron:23arxiv}, and Gemini~\citep{GeminiTeam:23arxiv,Reid:24arxiv}.

\section{Experiments}
\label{sec:Exp}

This section presents an empirical assessment of xLSTM, contrasting its performance with established approaches in the context of language modeling. Section~\ref{sec:ExpSyntethic} focuses on analyzing xLSTM's proficiency using artificially generated tasks. In Section~\ref{sec:ExpComparison}, we provide a comparative analysis of the validation set perplexities for a range of state-of-the-art language modeling approaches, all trained with 15B tokens from the SlimPajama corpus~\citep{Soboleva:23}. xLSTM-specific ablation studies are also performed on this dataset. Furthermore, we investigate the scaling properties of the various models in a manner similar to prior work by \citet{Kaplan:20} and \citet{Brown:20short}.

Section~\ref{sec:ExpLanguage} extends the investigation with a comprehensive language modeling study. Here, we examine xLSTM and top-performing baselines from Section~\ref{sec:ExpComparison}, now scaled up to 300B tokens of SlimPajama~\citep{Soboleva:23}. The analysis consists of four parts: (1) evaluating extrapolation abilities to longer input contexts, (2) assessing model performance based on validation perplexity and results on a suite of downstream tasks~\citep{Sutawika:24}, (3) testing across 571 text domains within the PALOMA language benchmark dataset~\citep{Magnusson:23arxivshort}, and (4) reassessing the models' scaling characteristics under a data regime 20 times larger than previously examined.

\label{sec:xlstm_blocks_notation}
Throughout all experiments, we denote the ratio of mLSTM to sLSTM xLSTM blocks as xLSTM[$a$:$b$], where $a/b$ indicates the count of mLSTM-based to sLSTM-based blocks, respectively. As an illustration, xLSTM[7:1] signifies that among eight blocks, seven utilize mLSTM and one uses sLSTM. When employing a total of 48 blocks, this configuration yields 42 mLSTM-based and 6 sLSTM-based blocks. Additionally, in every experiment, pre up-projection blocks are incorporated for mLSTM components, while post up-projection blocks are used for sLSTM components.

\subsection{Synthetic Tasks and Long Range Arena}
\label{sec:ExpSyntethic}
We begin by evaluating xLSTM’s novel exponential gating mechanism combined with memory mixing on a suite of formal language benchmarks~\citep{Deletang:23}. Next, we investigate the impact of the newly introduced matrix memory in xLSTM on the Multi-Query Associative Recall challenge~\citep{Arora:23arxiv}. Lastly, the capability of xLSTM to handle long input sequences is benchmarked using the Long Range Arena suite~\citep{Tay:21}.

\paragraph{Evaluation of xLSTM’s Exponential Gating with Memory Mixing.} We evaluate the proposed exponential gating mechanism with memory mixing in xLSTM, hypothesized to enhance its capabilities for state tracking tasks~\citep{Merrill:24, Merrill:23}. The formal language experiments introduced by \citet{Deletang:23} are implemented and expanded to support variable-length sequences for the purpose of testing length generalization. Detailed task descriptions and comprehensive results are provided in Appendix~\ref{sec:appExpSynthetic-formalLang}. xLSTM’s performance is benchmarked against a range of baseline architectures, including Transformers, State Space Models, and Recurrent Neural Networks. Performance is quantified by measuring accuracy exclusively on task-relevant tokens, with scores normalized between 0 (chance) and 1 (perfect). Specifically, we assess 2-block variants from multiple architectures: xLSTM[0:1] (pure sLSTM), xLSTM[1:0] (pure mLSTM), xLSTM[1:1], Llama, Mamba, RWKV, Retention, Hyena, standard LSTM, and LSTM as used within Transformer blocks (LSTM (Block)). The comparative findings are depicted in Figure~\ref{fig:formal-main}. Architectures such as Transformers and State Space Models that lack memory mixing and therefore do not maintain state tracking, are incapable of handling tasks based on regular grammars, such as the parity task. This is consistent with prior observations that Transformers and State Space models are inherently less expressive than RNN-based methods~\citep{Merrill:24, Merrill:23, Deletang:23}.

\paragraph{Test of xLSTM's Memory Capacities on Associative Recall Tasks.} This study evaluates the matrix-based memory of xLSTM in terms of its memory capacity using the Multi-Query Associative Recall benchmark~\citep{Arora:23arxiv}. In this task, each sequence is composed of randomly assigned key-value pairs selected from an extensive vocabulary; these pairs must be retained in memory for subsequent retrieval. To increase the complexity beyond that of the original evaluation, we scale up the number of key-value associations to as many as 256, and we expand the input context window to 2048 tokens. This setup allows for a more comprehensive examination of memory capacities across multiple models. We benchmark 2-block variants of Llama, Mamba, RWKV-5, RWKV-6, xLSTM[1:1], and xLSTM[1:0], with retrieval accuracy serving as the performance metric. Given that Transformer-based architectures (such as Llama) possess memory capacity that grows exponentially with coding dimension~\citep{Ramsauer:21}, they serve as the reference standard in this assessment. The primary results are summarized in Figure~\ref{fig:mqar-main}. Notably, xLSTM[1:1] outperforms other non-Transformer competitors, demonstrating strong performance even in models of smaller scale. It is noteworthy that incorporating the sLSTM block not only preserves but seems to further augment memory capacity, especially in the most challenging scenario involving 256 key-value pairs. Additional findings can be found in Appendix~\ref{sec:appExpSynthetic-mqar}, where analyses reveal that the improved memory mechanisms in xLSTM enable extrapolation to longer contexts than those encountered during training.

\paragraph{Test of xLSTM's Long Context Capabilities on Long Range Arena.} In order to evaluate xLSTM's ability to process extended sequences and maintain long-range dependencies, we conduct a comparative analysis of various models on the Long Range Arena~\citep{Tay:21}. xLSTM consistently achieves high accuracy across all evaluated tasks, indicating that the architecture excels at efficiently managing the challenges presented by long context scenarios.

Additional information is provided in Appendix~\ref{sec:appExpSynthetic-lra}.

\subsection{Method Comparison and Ablation Study}
\label{sec:ExpComparison}

In order to explore our central research question—namely, the potential of our novel LSTM variants when substantially scaled for language modeling tasks—we conduct experiments by training xLSTMs, Transformers, State Space Models, and several alternative approaches on a dataset comprising 15B tokens sourced from SlimPajama, all within an auto-regressive framework. The resulting models are evaluated using a shared validation set, and for xLSTMs, we undertake a series of ablation studies to further dissect their performance.

\paragraph{Comparison of xLSTM with Existing Approaches.} To facilitate a meaningful comparison, all models are trained using a corpus of 15B tokens from SlimPajama~\citep{Soboleva:23}. We assess model performance based on their perplexity evaluated on a held-out validation set. The methods included in this comparison are: our proposed xLSTM, GPT-3 (Transformer)~\citep{Brown:20short}, Llama (Transformer)~\citep{Touvron:23arxiv}, H3 (SSM)~\citep{Fu:23}, Mamba (SSM)~\citep{Gu:24arxiv}, RWKV-4 (RNN)~\citep{Peng:23arxivshort}, RWKV-5 (RNN)~\citep{Peng:24arxivshort}, RWKV-6 (RNN)~\citep{Peng:24arxivshort}, GLA (linear Transformer)~\citep{Yang:23arxiv}, HGRN (RNN)~\citep{Qin:23}, HGRN2 (RNN)~\citep{Qin:24arxiv}, RetNet (linear Transformer)~\citep{Sun:23arxiv}, Hyena (linear Transformer)~\citep{Poli:23}, as well as xLSTM[1:0] and xLSTM[7:1] variants (refer to Section~\ref{sec:xlstm_blocks_notation}). Training was carried out using mixed precision; however, for RWKV-5, RWKV-6, GLA, and HGRN2, mixed precision training was performed without relying on PyTorch’s automatic mixed precision (see also Appendix Section~\ref{sec:appGenTrainProc} for specifics).

We divide the approaches into three groups: (a) Transformers, (b) State Space Models (SSMs), and (c) Recurrent Neural Networks (RNNs) and linear Transformers. Linear Transformers comprise methods that employ a linear alternative to standard attention. All model architectures are matched to a GPT-3 configuration containing 350M parameters, employing a 1024-dimensional embedding and 24 residual blocks. GPT-3 is unique in that it shares token and output embedding weights, resulting in a slightly lower parameter count. Results reported in Table~\ref{tab:spaj15b_model_results} demonstrate that xLSTM achieves the lowest validation perplexity among all models compared. Additional information is available in Appendix~\ref{sec:appExpComparison}. Furthermore, Figure~\ref{fig:temp_sclaw15B} depicts scaling trends for these experiments, suggesting that the performance advantages of xLSTM persist as model size increases.

\paragraph{Ablation Studies.}
As shown in Table~\ref{tab:spaj15b_model_results} and Figure~\ref{fig:temp_sclaw15B}, the xLSTM model demonstrates superior performance in language modeling tasks when trained on 15B tokens from the SlimPajama dataset. In order to isolate the individual contributions that distinguish xLSTM from standard LSTM, we progressively modify a baseline LSTM architecture into an xLSTM through a series of controlled ablations. The first modification involves inserting LSTM layers within pre-LayerNorm residual backbones. In the second step, we incorporate a post up-projection block. The final transformation consists of adding exponential gating mechanisms along with matrix memory capabilities.

Table~\ref{tab:ablstudies} (top) presents these findings.

The ablation experiments reveal that the notable gains in performance arise from the integration of both exponential gating and matrix memory. Recognizing the critical role gating plays within RNNs and State Space Models, we further examine various gating strategies. As indicated in Table~\ref{tab:ablstudies} (bottom), enabling each gate to be both learnable and directly affected by the input consistently improves results. Further details regarding the ablation of specific backbone elements are provided in Appendix~\ref{sec:appAblDetails}.

\subsection{xLSTM as Large Language Model}
\label{sec:ExpLanguage}

This investigation concludes with comprehensive language modeling experiments at scale, evaluating xLSTM's capabilities as a large language model (LLM). To this end, we expand the training dataset and utilize 300B tokens drawn from SlimPajama. Notably, an equivalent volume of training data is adopted by other models such as Mamba~\citep{Gu:24arxiv} and Griffin~\citep{De:24arxiv}.

Our study benchmarks xLSTM against RWKV-4, Llama, and Mamba, each serving as a representative of their respective methodological categories, as outlined in Section~\ref{sec:ExpComparison}. RWKV-4 was chosen as the RNN baseline because appropriate training precision settings for RWKV-5, RWKV-6, and HGRN2 were only established following the initiation of the 300B token training runs (refer to Appendix~\ref{sec:appGenTrainProc}). We experiment with several model scales (125M, 350M, 760M, 1.3B), evaluating all models for their ability to generalize to longer sequences, and measuring their accuracy on the validation dataset. Further, we examine downstream task performance, conduct comprehensive language modeling assessments using the 571 domains from the PALOMA benchmark, and analyze the models’ scaling law characteristics.

\paragraph{Sequence Length Extrapolation.}
\label{sec:spaj300B_ctx_extrapolation}
We begin by evaluating the ability of large 1.3B-parameter models—namely xLSTM, RWKV-4, Llama, and Mamba—to extrapolate to longer sequence lengths. Each model is trained with a context window of 2048 tokens and subsequently assessed on sequences extending up to 16384 tokens. Refer to Figure~\ref{fig:spaj300B_seq_extrapolation} for a visual comparison of performance. Notably, xLSTM achieves consistently low perplexity as the input length increases, outperforming the alternative approaches at larger context sizes.

\paragraph{Validation Perplexity and Downstream Tasks.}
\label{sec:spaj_downstream_eval}
In the second part of our analysis, we assess all model variants—including xLSTM, RWKV-4, Llama, and Mamba—by measuring their perplexity on the SlimPajama validation dataset for next-token prediction, as well as their abilities on downstream benchmarks focused on evaluating common sense reasoning.

The third column in Table~\ref{tab:lmeval} presents the validation set perplexity scores for various approaches. Across all model capacities, both xLSTM[1:0] and xLSTM[7:1] demonstrate superior performance in terms of validation perplexity. The remaining columns in Table~\ref{tab:lmeval} summarize results on downstream tasks, where xLSTM consistently outperforms other methods for nearly all tasks and model scales; Mamba achieves the top result exclusively on the ARC task in select configurations. Further information is available in Appendix~\ref{sec:appExpLanguage}.

\paragraph{Performance on PALOMA Language Tasks.} Thirdly, we evaluate xLSTM, RWKV-4, Llama, and Mamba models—across all parameter scales—on next token prediction for the PALOMA language task suite~\citep{Magnusson:23arxivshort}. Specifically, model performance is assessed via perplexity scores computed across 571 distinct textual domains, spanning sources such as nytimes.com and Reddit’s r/depression. Table~\ref{tab:ppleval} presents the perplexity results, organizing them into language modeling tasks (first seven columns) and specialized domain-level benchmarks (final five columns). Among the xLSTM configurations, xLSTM[1:0] yields superior results compared to xLSTM[7:1] for these tasks.

xLSTM[1:0] demonstrates superior perplexity performance compared to Mamba across 568 of 571 (99.5\%) text domains, outperforms Llama in 486 out of 571 (85.1\%) domains, and surpasses RWKV-4 in 570 out of 571 (99.8\%) domains. Additional information is provided in Appendix~\ref{sec:appExpLanguage}.

\paragraph{Scaling Laws.} Fourth, we examine the power-law scaling properties, which facilitate predicting model performance as parameter count increases~\citep{Kaplan:20,Brown:20short}. Figure~\ref{fig:temp_sclaw300B} illustrates this scaling phenomenon. The evaluated architectures demonstrate comparable scaling trends, though they exhibit distinct baseline performance levels. Among these, RWKV-4 achieves the lowest results, followed in sequence by Llama and Mamba. xLSTM surpasses Mamba by a margin similar to that observed between Mamba and Llama. This scaling pattern suggests that as models are scaled up, xLSTM is expected to maintain superior performance relative to both Transformer-based and State-Space counterparts.

\textbf{Generation Times and Maximal Throughput.} We present our measurements of text generation latency in Figure~\ref{fig:inference_time_generation} and report the peak decoding throughput in Figure~\ref{fig:inference_time_decoding_throughput} (left) for the various xLSTM models at the 1.3B parameter scale. Our analysis includes comparisons with equivalently sized versions of Mamba, Llama, and RWKV from HuggingFace repositories, noting that our Llama baseline leverages a static key-value cache. During these experiments, neither the compilation of the entire cache for the Transformer model nor the compilation of the Mamba implementation with \texttt{torch.compile} was successful. For all text generation tests, we utilize a batch size of 1 and set the pre-fill to 16, favoring the Transformer architecture. 
In Figure~\ref{fig:inference_time_generation}, one can observe that both xLSTM and other recurrent models like Mamba and RWKV-4 scale their generation times linearly, whereas Llama's latency increases quadratically. For the assessment of decoding throughput, a range of batch sizes and prefill settings were explored for Llama. 
As illustrated in Figure~\ref{fig:inference_time_decoding_throughput} (right), xLSTM supports substantially larger batch sizes relative to Llama thanks to its fixed memory requirements, resulting in superior throughput performance.

\section{Limitations}

(i) Unlike mLSTM, the design of memory mixing in sLSTM restricts opportunities for parallel execution, thereby preventing rapid parallel implementations. Despite this, we created an efficient CUDA kernel for sLSTM, which at present operates at less than double the execution time compared to our parallelized mLSTM implementation. (ii) The mLSTM CUDA kernels currently lack optimization; as a result, their execution speed is approximately four times slower than what is achieved with FlashAttention or the scan in Mamba. There is potential for substantial speedup through kernel optimizations similar to FlashAttention. (iii) The use of $d \times d$ matrices in mLSTM memory introduces significant computational overhead. However, since memory update and retrieval are parameter-free operations that lend themselves to parallelization through typical matrix operations, the added computational time from memory complexity is minimal in practice. (iv) Careful selection of the forget gate initialization values is necessary. (v) Because the matrix memory’s capacity is fixed with respect to sequence length, extending sequence length for larger contexts may saturate the memory. Nonetheless, this issue has not manifested for context lengths up to 16k tokens, as discussed in Section~\ref{sec:spaj300B_ctx_extrapolation}. (vi) Owing to the high computational cost of training large language models, neither the architecture nor the hyperparameters were thoroughly tuned, particularly for larger xLSTM variants. It is expected that substantial architectural and hyperparameter optimization will be required for xLSTM models to realize their full capabilities.

\section{Conclusion}
Our investigation has provided a partial response to the question: How effective is language modeling when LSTM architectures are scaled to billions of parameters? The evidence so far demonstrates that they can achieve performance comparable to prevailing models, such as Transformers and State Space Models. By introducing xLSTM—featuring exponential gating combined with memory mixing and an updated memory configuration—we present a variant that matches or surpasses state-of-the-art approaches in language modeling tasks. Observed scaling laws suggest that expanding xLSTM models may pose a strong challenge to today’s Large Language Models reliant on Transformer architectures. Furthermore, xLSTM shows significant promise for application areas beyond language modeling, including Reinforcement Learning, Time Series Forecasting, and physical systems modeling.

\end{document}